\section*{الفصل \ref{ch:g9:periodic-table-first-look} --- الجدول الدوري: نظرة أولى}

\begin{solution}{exo:g9:periodic-table-first-look:1}
بحسب تزايد \omterm{def:g9:inside-the-atom:atomic-number}{العدد الذري} $Z$.
\end{solution}

\begin{solution}{exo:g9:periodic-table-first-look:2}
الكربون: \omterm{def:g9:periodic-table-first-look:period}{الدورة} 2، \omterm{def:g9:periodic-table-first-look:period}{الزمرة} 14. المغنيسيوم: \omterm{def:g9:periodic-table-first-look:period}{الدورة} 3، \omterm{def:g9:periodic-table-first-look:period}{الزمرة} 2. الأرغون:
\omterm{def:g9:periodic-table-first-look:period}{الدورة} 3، \omterm{def:g9:periodic-table-first-look:period}{الزمرة} 18.
\end{solution}

\begin{solution}{exo:g9:periodic-table-first-look:3}
\omterm{def:g9:periodic-table-first-look:metal}{فلزات}: الحديد، والنحاس، والألومنيوم. لافلزات: الكبريت، والأكسجين، والكلور.
\end{solution}

\begin{solution}{exo:g9:periodic-table-first-look:4}
البوتاسيوم (\ce{K}) تحت الصوديوم؛ والبروم (\ce{Br}) تحت الكلور.
\end{solution}

\begin{solution}{exo:g9:periodic-table-first-look:5}
الكالسيوم (\ce{Ca}، $Z = 20$) والنيتروجين (\ce{N}، $Z = 7$).
\end{solution}

\begin{solution}{exo:g9:periodic-table-first-look:6}
فهو في \omterm{def:g9:periodic-table-first-look:period}{الزمرة} نفسها التي فيها الصوديوم، \omterm{def:g9:periodic-table-first-look:period}{الزمرة} 1، وعناصر \omterm{def:g9:periodic-table-first-look:period}{الزمرة} الواحدة
تتشابه في سلوكها.
\end{solution}

\begin{solution}{exo:g9:periodic-table-first-look:7}
\omterm{def:g9:periodic-table-first-look:metal}{فلز}: فهو لامع، ويوصل الكهرباء، ويكوّن أيونًا موجبًا.
\end{solution}

\begin{solution}{exo:g9:periodic-table-first-look:8}
ليُبقي العناصر المتشابهة الخصائص في العمود نفسه، حتى لو اقتضى ذلك ترك خانة
فارغة. وكانت الفجوات تنبؤات: فقد وُجدت العناصر المفقودة لاحقًا بالخصائص التي
وصفها.
\end{solution}

\begin{solution}{exo:g9:periodic-table-first-look:9}
الكربون: $\ce{C} = 12$؛ والأكسجين: $\ce{O} = 16$. والفجوتان «? $= 68$» (بجوار
Al $= 27.4$) و«? $= 70$» (بجوار Si $= 28$).
\end{solution}

\begin{solution}{exo:g9:periodic-table-first-look:10}
\omterm{def:g9:periodic-table-first-look:period}{الدورة} 2: 8 عناصر (من الليثيوم إلى النيون). \omterm{def:g9:periodic-table-first-look:period}{الدورة} 4: 18 عنصرًا (من البوتاسيوم
إلى الكريبتون).
\end{solution}

\begin{solution}{exo:g9:periodic-table-first-look:11}
\omterm{def:g9:periodic-table-first-look:period}{الزمرة} 18. فيما يجب ألا يتفاعل فيه شيء: الأرغون داخل المصابيح وفي غاز
اللحام، والهيليوم في البالونات (فهو لا يحترق).
\end{solution}

\begin{solution}{exo:g9:periodic-table-first-look:12}
$(27.0 + 114.8)/2 = 70.9$، وهو قريب من 69.7.
\end{solution}

\begin{solution}{pb:g9:periodic-table-first-look:1}
\textbf{1.} $Z = 32$؛ \omterm{def:g9:periodic-table-first-look:period}{الدورة} 4، \omterm{def:g9:periodic-table-first-look:period}{الزمرة} 14.

\textbf{2.} فوقه: السيليكون. تحته: القصدير. عن يساره: الغاليوم. عن يمينه: الزرنيخ.

\textbf{3.} شبه \omterm{def:g9:periodic-table-first-look:metal}{فلز}.

\textbf{4.} فهو في \omterm{def:g9:periodic-table-first-look:period}{الزمرة} نفسها التي فيها السيليكون، وعناصر \omterm{def:g9:periodic-table-first-look:period}{الزمرة} الواحدة
تتشابه في سلوكها.

\textbf{5.} $(28.1 + 118.7)/2 = 73.4$.

\textbf{6.} $(69.7 + 74.9)/2 = 72.3$.

\textbf{7.} متوسط اليسار واليمين (72.3): فعلى امتداد \omterm{def:g9:periodic-table-first-look:period}{الدورة} يزداد الوزن الذري
بانتظام، عنصرًا بعد عنصر، أما في النزول عبر \omterm{def:g9:periodic-table-first-look:period}{الزمرة} فيقفز قفزات كبيرة.

\textbf{8.} $72.6 - 70 = 2.6$.

\textbf{9.} $(28.1 + 118.7 + 69.7 + 74.9)/4 = 291.4/4 = 72.85$، أي
72.9.

\textbf{10.} المتوسط، 72.9، قريب من 72 عند مندلييف ومن 72.3 عند فينكلر
(ومن 72.6 اليوم).

\textbf{11.} عنصر تُنبئ به قبل أن يراه أحد وُجد بالخصائص المتوقعة: فلم يكن
الجدول مجرد قائمة، بل كان قادرًا على التنبؤ.

\textbf{12.} نحو 72.9.
\end{solution}
